\item Use these equations to show that in this economy
\begin{equation}\begin{gathered}\begin{aligned}
        \dot{C} & =  (\rfree-\timeRate) C - \pDies(\pDies+\timeRate) W \label{eq:aggCdot}
\end{aligned}\end{gathered}\end{equation}

Hint: Differentiate \eqref{eq:clevel} and substitute out for $H$ by solving
\eqref{eq:clevel} for $H$.

\ifthenelse{\boolean{AnswersTF}}{

\begin{quote}
{\it Answer:}

Time differentiate \eqref{eq:clevel} and substitute for $\dot{H}$, 
$\dot{W}$, and $H$ to get
\begin{equation*}\begin{gathered}\begin{aligned}
        \dot{C} & =  (\pDies+\timeRate)\left[(\rfree+\pDies) H - Y + \rfree W + Y - C\right]  
\\       & =  (\pDies+\timeRate)\left[(\rfree+\pDies) H + \rfree W \right] - C (\pDies+\timeRate)
\\       & =  (\pDies+\timeRate)\left[(\rfree+\pDies) \left(\frac{C}{\pDies+\timeRate}-W\right)+r W\right] - C (\pDies+\timeRate)
\\   & =  (\rfree+\pDies) C + (\pDies+\timeRate)\left[    rW - (\rfree+\pDies) W \right] - C(\pDies+\timeRate)
\\   & =  (\rfree-\timeRate) C - (\pDies+\timeRate)\pDies W.
\end{aligned}\end{gathered}\end{equation*}

\end{quote}
}{\pagebreak

.
\pagebreak
}

Now assume there is a standard production function
$\FFunc(K)=K^{\alpha}-\delta K$ and assume perfect competition so that the
net interest rate $r$ is equal to the net marginal product of capital,
\begin{equation}
        \rfree = \FFunc^{\prime}(K) = \alpha K^{\alpha-1}-\delta
\end{equation}
and the aggregate capital stock at time $t$ is the same as aggregate
nonhuman wealth, $K(t)=W(t)$.  The aggregate accumulation equation is
just the usual
\begin{equation}\begin{gathered}\begin{aligned}
  \dot{K} & =  K^{\alpha}-\delta K -C. \label{eq:aggkdot}
\end{aligned}\end{gathered}\end{equation}

\item Show that the following equation 
\begin{equation}\begin{gathered}\begin{aligned}
  \pDies(\pDies+\timeRate) K      & =  C(\alpha K^{\alpha-1}-\delta-\timeRate)  \label{eq:CvsK1}
\end{aligned}\end{gathered}\end{equation}
describes the $\dot{C}=0$ locus.  Use this equation to show that
\begin{equation}\begin{gathered}\begin{aligned}
  \label{eq:1}
  \lim_{C \rightarrow 0} K & =  0
\\ \lim_{C \rightarrow \infty} K & =  ((\timeRate+\delta)/\alpha)^{1/{\alpha-1}}.
\end{aligned}\end{gathered}\end{equation}

\ifthenelse{\boolean{AnswersTF}}{
\begin{quote}
{\it Answer:}

Rewriting the $\dot{C}$ equation as a function of $K$ yields
\begin{equation}\begin{gathered}\begin{aligned}
  \dot{C} & =  (\alpha K^{\alpha-1} - \timeRate - \delta)C - (\pDies+\timeRate)\pDies K.
\end{aligned}\end{gathered}\end{equation}

The $\dot{C}=0$ locus is therefore given by
\begin{equation}\begin{gathered}\begin{aligned}
      0 & =  (\alpha K^{\alpha-1}-\timeRate-\delta) - \pDies(\pDies+\timeRate) K/C \label{eq:aggCdotmod}
\\  \pDies(\pDies+\timeRate) K/C       & =  (\alpha K^{\alpha-1}-\delta-\timeRate) 
\\  \pDies(\pDies+\timeRate) K      & =  C(\alpha K^{\alpha-1}-\delta-\timeRate)  
\end{aligned}\end{gathered}\end{equation}

In the limit as $C \rightarrow 0$ this expression approaches
\begin{equation}
  \label{eq:2}
  \pDies(\pDies+\timeRate) K = 0
\end{equation}
which can be true only if $\lim_{\{C \rightarrow 0\}} K = 0$.  

On the other hand, as $C \rightarrow \infty$ we know that $K$
must remain finite (because the DBC does not allow infinite
accumulation of $K$ -- infinite $K$ would imply infinite depreciation which 
could never be paid for by a production function with diminishing marginal returns), which means that 
\begin{equation}\begin{gathered}\begin{aligned}
\lim_{\{C \rightarrow \infty\}} \alpha K^{\alpha-1} & =  (\timeRate+\delta)
\\ K^{*} & =  ((\timeRate+\delta)/\alpha)^{1/{\alpha-1}}.
\end{aligned}\end{gathered}\end{equation}

\end{quote}
}{\pagebreak

}

\item Using your results from the previous question and the aggregate
  dynamic budget constraint, draw the phase diagram for this model and
  discuss how and why it differs from the phase diagram for the
  standard Ramsey-Cass/Koopmans model.  (please use $K^{*}$ to
  designate the RCK steady-state capital stock and $\bar{K}$ to
  designate this model's steady-state $K$.)

\ifthenelse{\boolean{AnswersTF}}{
\begin{quote}
{\it Answer:}

Using \eqref{eq:aggkdot}, the $\dot{K}=0$ locus is
\begin{equation}\begin{gathered}\begin{aligned}
  C & =  K^{\alpha}-\delta K
\end{aligned}\end{gathered}\end{equation}
which yields the usual hump-shaped $\dot{K}=0$ locus.

Rewriting the $\dot{C}$ equation as a function of $K$ yields
\begin{equation}\begin{gathered}\begin{aligned}
  \dot{C} & =  (\alpha K^{\alpha-1} - \timeRate - \delta)C - (\pDies+\timeRate)\pDies K.
\end{aligned}\end{gathered}\end{equation}

The $\dot{C}=0$ locus is therefore given by
\begin{equation}\begin{gathered}\begin{aligned}
      0 & =  (\alpha K^{\alpha-1}-\timeRate-\delta) - \pDies(\pDies+\timeRate) K/C 
\\  \pDies(\pDies+\timeRate) K/C       & =  (\alpha K^{\alpha-1}-\delta-\timeRate) 
\\  \pDies(\pDies+\timeRate) K      & =  C(\alpha K^{\alpha-1}-\delta-\timeRate)  \label{eq:CvsK}
\end{aligned}\end{gathered}\end{equation}

In the limit as $C \rightarrow 0$ this expression approaches
\begin{equation}
  \pDies(\pDies+\timeRate) K = 0
\end{equation}
which can be true only if $\lim_{\{C \rightarrow 0\}} K = 0$.  

On the other hand, as $C \rightarrow \infty$ we know that $K$
must remain finite (because the DBC does not allow infinite
accumulation of $K$), which means that 
\begin{equation}\begin{gathered}\begin{aligned}
\lim_{\{C \rightarrow \infty\}} \alpha K^{\alpha-1} & =  (\timeRate+\delta)
\\ K^{*} & =  ((\timeRate+\delta)/\alpha)^{1/{(\alpha-1)}}.
\end{aligned}\end{gathered}\end{equation}

\begin{figure}
\centerline{\includegraphics[width=6in]{../Figures/BlanchardFiniteHorizPhase1}}
\end{figure}

In the infinite horizon economy we have $\pDies = 0$ and so the
steady-state interest rate would be the $K^{*}$ where $\alpha
K^{\alpha-1}-\delta = \timeRate$.  But since $K/C$ and $\pDies(\pDies+\timeRate)$ are
strictly positive, in this finite-horizon economy we would have
$\dot{C}/C < 0$ at $K=K^{*}$.  It is clear therefore that in order
for \eqref{eq:aggCdotmod} to hold we will need $\alpha K^{\alpha-1}$
to be larger than it is at $K^{*}$, which is to say we need a higher
steady-state interest rate, and thus we need a lower steady-state
capital stock, which is depicted in the figure as $\bar{K}$.

This makes sense because the finite-horizon consumers in this economy 
discount the future more than the representative agent does, because 
they die but a representative agent does not.

These results are combined in the figure, which shows that the intersection
of the $\dot{C}=0$ locus intersects the $\dot{K}=0$ locus at point $A$ 
which corresponds to a lower level of the capital stock than in the 
infinite horizon model.  

\end{quote}
}{\pagebreak

}

Now consider the introduction of a government that finances spending
either by lump-sum taxes or by debt.  Its dynamic budget constraint is
\begin{equation}
  \dot{D} =  \rfree D + G - \TaxLev
\end{equation}
where $D$ is government debt, $G$ is government spending, and
$\TaxLev$ is a lump-sum per capita tax.  Defining
\begin{equation}\begin{gathered}\begin{aligned}
  \mathcal{R}_t^{s} = e^{\int_t^{s}\rfree_v dv}
\end{aligned}\end{gathered}\end{equation}
as the compound interest factor between time $t$ and time $s$, the
government is also required to satisfy the transversality condition
\begin{equation}
  \label{eq:10}
  \lim_{t \rightarrow \infty} D_t/\mathcal{R}_t^{s} = 0.
\end{equation}

Consider the following fiscal policy experiment.  Until time $t$ there
has been no government ($G_s = D_s = \TaxLev_s=0~\forall~s<t$).  At date
$t$ the government issues a quantity $D$ of debt and announces that
future lump sum taxes will be imposed in amounts exactly large enough
to pay the interest on this debt (so subsequently, $\dot{D}=0$
forever).  The government rebates the proceeds of its sale of debt to
the public as a per-capita lump sum of $D$ per person.  The government
will never engage in any spending (aside from paying interest on the
debt).  Define the new variables
\begin{equation}\begin{gathered}\begin{aligned}
  \mathcal{W} & =  K + D
\\ \mathcal{Y} & =  Y - \TaxLev
\\ \mathcal{H} & =  \int_{t}^{\infty} \mathcal{Y}/\hat{\mathcal{R}}_{t}^{s} ds
\end{aligned}\end{gathered}\end{equation}

\item Explain why the effect of this policy is to modify the aggregate
specification of the economy to (see top of next page) \pagebreak
\begin{equation}\begin{gathered}\begin{aligned}
   C & =  (\pDies+\timeRate) (\mathcal{H} + \mathcal{W})  \label{eq:Cnew}
\\ \dot{\mathcal{W}} & =  \rfree \mathcal{W} + \mathcal{Y}- C 
\\ \dot{\mathcal{H}} & =  (\rfree+\pDies) \mathcal{H} - \mathcal{Y}
\end{aligned}\end{gathered}\end{equation}

\ifthenelse{\boolean{AnswersTF}}{\begin{quote}{\it Answer:}
    
    The effect of the government policy is twofold.  On the one hand,
    the distribution of government bonds increases the consumers'
    wealth $W$ by an amount equal to the value of the bonds received,
    resulting in a new definition of wealth $\mathcal{W}$ which
    includes the bonds.  On the other hand, the higher value of taxes
    off to infinity reduces the consumers' human wealth by an amount
    equal to the present discounted value of the taxes.  
    
    The change in (redefined) wealth is now {\it net} income $Y-\timeRate$
    minus consumption.


\end{quote}}{\pagebreak }

\item Use these equations to show that the new dynamic equations for
the economy are
\begin{equation}\begin{gathered}\begin{aligned}
  \dot{C} & =  (\rfree - \timeRate)C - \pDies(\pDies+\timeRate)(K+D)
\\ \dot{K} & =  K^{\alpha}-\delta K - C
\end{aligned}\end{gathered}\end{equation}

Show how the policy change affects the economy over time, using a
phase diagram and a diagram showing the dynamics of aggregate
consumption after the policy is introduced.  Explain the impact
of the policy on different generations in the economy.

\ifthenelse{\boolean{AnswersTF}}{\begin{quote}{\it Answer:}

Time differentiating \eqref{eq:Cnew} yields
\begin{equation}\begin{gathered}\begin{aligned}
  \dot{C} & =  (\pDies+\timeRate) (\dot{\mathcal{H}} + \dot{\mathcal{W}})
\\ & =  (\pDies+\timeRate)(\rfree \mathcal{W}+(\rfree+\pDies)\mathcal{H}-C) \label{eq:dotCnew}
\end{aligned}\end{gathered}\end{equation}

Now solve \eqref{eq:Cnew} for $\mathcal{H}$,
\begin{equation}\begin{gathered}\begin{aligned}
  \label{eq:18}
  \mathcal{H} & =  \left(\frac{C}{\pDies+ \timeRate}\right)-\mathcal{W}
\end{aligned}\end{gathered}\end{equation}
and substitute into \eqref{eq:dotCnew} to obtain
\begin{equation}\begin{gathered}\begin{aligned}
  \dot{C} & =  (\pDies+\timeRate)(\rfree \mathcal{W}+(\rfree+\pDies)(C(\pDies+\timeRate)^{-1} - \mathcal{W})-C) 
\\        & =  (\rfree+\pDies)C - (\pDies+\timeRate)C- \pDies(\pDies+\timeRate)\mathcal{W}
\\ & =  (\rfree-\timeRate)C - \pDies(\pDies+\timeRate)(K+D).
\end{aligned}\end{gathered}\end{equation}

Since we are assuming that $D$ is a constant (after the fiscal
experiment), any combination of $C$ and $W$ that would have been on
the $\dot{C}=0$ locus before the policy shift now has a value $\dot{C}
= -\pDies(\pDies+\timeRate)D$.  This means that the $C$ that would restore
$\dot{C}=0$ must be a larger $C$, which says that the $\dot{C}=0$
locus shifts up (or, equivalently, to the left).  Thus, the new
equilibrium will be at a lower value of $K$ and a higher interest
rate.

The phase diagram shows that the new equilibrium point $A'$ is to the left
of the original equilibrium.  This is because at a given level of the aggregate
capital stock, consumers spend more because $\mathcal{W}>W$.  Thus, 
Ricardian equivalence does not hold in this model, because a tax cut today
financed by a future perpetual tax is a transfer of resources from future 
consumers to today's consumers, and there are no altruistic links that
make current consumers offset this by saving more on behalf of future
generations.

\begin{figure}
\centerline{\includegraphics[width=6in]{../Figures/BlanchardFiniteHorizPhase2}}
\end{figure}

The next figure shows the path of consumption per capita in this
economy.  Prior to time 0, the economy was in its steady-state
equilibrium at the level of consumption $C_0$ corresponding to the
equilibrium labeled $A$ in the phase diagram.  At time 0, the fiscal
policy is carried out.  The fiscal policy immediately increases
consumption because it amounts to a transfer of resources from future
to current generations.  However, the higher level of consumption runs
down the capital stock per capita, and so over time consumption asypmtotically
approaches a new, lower equilibrium level of consumption $C'$.

\begin{figure}
\centerline{\includegraphics[width=6in]{../Figures/BlanchardCDynamics}}
\end{figure}


\end{quote}
}{\pagebreak}

\item Blanchard shows that if the model is changed so that each
  consumer's income declines exponentially at rate $\gamma$ after
  birth, the result is equivalent to assuming that future labor income
  is discounted at an interest rate that is higher by $\gamma$.  He
  further shows that the equations of motion of the model change to
\begin{equation}\begin{gathered}\begin{aligned}
  \dot{C} & =  (\alpha K^{\alpha-1}-\delta +\gamma-\timeRate)C-(\pDies+\gamma)(\pDies+\timeRate)K \label{eq:lccdot}
\\ \dot{K} & =  K^{\alpha}-\delta K - C. %\label{eq:lckdot}
\end{aligned}\end{gathered}\end{equation}

Note that it is possible to rewrite the $\dot{C}=0$ locus as 
\begin{equation}\begin{gathered}\begin{aligned}
  (\alpha K^{\alpha-1}-\delta +\gamma-\timeRate)C & =  (\pDies+\gamma)(\pDies+\timeRate)K \label{eq:lccdoteq0}
\end{aligned}\end{gathered}\end{equation}

Use this equation to show that as $C \rightarrow \infty$, the
$\dot{C}=0$ locus asymptotes to
\begin{equation}\begin{gathered}\begin{aligned}
 \rfree = \alpha K^{\alpha-1} - \delta & =  \timeRate - \gamma \label{eq:eqblcr}
\end{aligned}\end{gathered}\end{equation}

Thus, for a large enough value of $\gamma$ it is possible that the net
interest rate in this economy could be negative.  Draw a phase diagram
corresponding to an equilibrium with a negative net interest rate, and
comment on why this is an interesting case to think about.  In
particular, what new light does it shed on the fiscal policy
experiment examined above?


\ifthenelse{\boolean{AnswersTF}}{\begin{quote}{\it Answer:}

As $C$ goes to infinity on the LHS of \eqref{eq:lccdoteq0}, the only way
the equation can continue to hold is if 
\begin{equation}\begin{gathered}\begin{aligned}
 \lim_{C \rightarrow \infty} (\alpha K^{\alpha-1}-\delta +\gamma-\timeRate) & =  0
\end{aligned}\end{gathered}\end{equation}
which implies \eqref{eq:eqblcr}. 

The new phase diagram shows the $\dot{C}=0$ locus intersecting the 
$\dot{K}=0$ locus to the right of the maximum of the $\dot{K}=0$ locus.  
This is implied by the fact that the net interest rate is negative, which
means that the net interest rate could be increased by reducing
the capital stock.

The reason this is an interesting case is that in this case it is 
possible for the economy to be in a condition of dynamic inefficiency,
just as in the 2-period OLG models discussed early in the class.  The idea
is to think of declining labor income as a way to generate a `life cycle' 
saving motive.  In such a case the fiscal experiment examined above
is interesting because it could rescue an economy with too much capital
from a state of dynamic inefficiency. 

\begin{figure}
\centerline{\includegraphics[width=6in]{../Figures/BlanchardFiniteHorizPhase3}}
\end{figure}

\end{quote}
}{\pagebreak}
